\documentclass[11pt]{article}
\usepackage[utf8]{inputenc}
\usepackage[T1]{fontenc}
\usepackage{lmodern}
\usepackage[spanish]{babel}
\usepackage{amsmath,amssymb}
\usepackage{graphicx}
\usepackage[left=2cm,right=2cm,top=2cm,bottom=2cm]{geometry}
\usepackage{enumitem}
\usepackage{booktabs}
\usepackage{array}
\usepackage{xcolor}
\usepackage{framed}

\spanishdecimal{.}

\setlist[enumerate,1]{label=\textbf{\arabic*.}, leftmargin=*, itemsep=6pt}
\setlist[enumerate,2]{label=(\alph*), itemsep=4pt}
\setlength{\parindent}{0pt}
\setlength{\parskip}{4pt}
\pagestyle{plain}

\newcommand{\Ein}{E_{\text{in}}}
\newcommand{\Eout}{E_{\text{out}}}
\newcommand{\Hyp}{\mathcal{H}}
\newcommand{\dvc}{d_{\text{VC}}}
\newcommand{\vx}{\mathbf{x}}
\newcommand{\vw}{\mathbf{w}}
\newcommand{\vy}{\mathbf{y}}
\newcommand{\R}{\mathbb{R}}
\newcommand{\T}{^{\top}}
\newcommand{\E}{\mathbb{E}}
\newcommand{\Prob}{\mathbb{P}}
\newcommand{\pts}[1]{\hfill\textbf{(#1~puntos)}}
\newcommand{\linea}[1]{\rule{#1}{0.4pt}}
\newcommand{\ptsinc}[1]{{\unskip\nobreak\hfil\penalty50\hskip1em\hbox{}\nobreak\hfill(#1)\parfillskip=0pt\finalhyphendemerits=0\par}}

\newcommand{\encabezado}[2]{%
\begin{center}
{\large\bfseries Reconocimiento de Patrones}\\[2pt]
Universidad de Sonora \textperiodcentered{} Licenciatura en Ciencias de la Computación\\
Semestre 2027-1 \textperiodcentered{} Prof.\ Julio Waissman Vilanova\\[10pt]
{\Large\bfseries #1}%
\ifdefined\tipodoc\\[4pt]{\large\itshape\tipodoc}\fi
\end{center}
\ifdefined\tipodoc\else
\vspace{6pt}
\textbf{Nombre:} \hrulefill\quad \textbf{Expediente:} \rule{3.5cm}{0.4pt}\par
\vspace{6pt}
\fi
#2
\par\medskip\hrule\medskip
}

\newcommand{\espacio}[1]{\par\vspace{#1}}
\newcommand{\nuevapagina}{\newpage}

\usepackage{pgfplots}
\pgfplotsset{compat=1.18}
\usepgfplotslibrary{groupplots}

\begin{document}

\encabezado{Examen rápido 3 \textperiodcentered{} Fundamentos del aprendizaje supervisado}{%
Tiempo: 20 minutos. Responde en esta hoja. Puedes usar calculadora; no se permiten apuntes.}

\begin{enumerate}

%% ---------------------------------------------------------------
\item \textbf{Discos centrados en el origen.} Sea $\Hyp$ el conjunto de clasificadores en $\R^2$
\[
h_r(\vx) = \begin{cases} +1 & \text{si } \lVert \vx \rVert \le r, \\ -1 & \text{si no,} \end{cases} \qquad r \ge 0 .
\]
\begin{enumerate}
  \item Para $N$ puntos con normas distintas, ¿cuántas dicotomías produce $\Hyp$? Es decir, ¿cuánto vale $m_{\Hyp}(N)$? Explica.
  \espacio{2.8cm}

  \item ¿Cuánto vale $\dvc(\Hyp)$? Los puntos viven en $\R^2$, ¿por qué no es 3, como el perceptrón en el plano?
  \espacio{2.5cm}

  \item Ahora permitimos también la orientación contraria: $\Hyp' = \{h_r\} \cup \{-h_r\}$ («dentro $+1$» o «dentro $-1$»). ¿Cuánto valen $m_{\Hyp'}(N)$ y $\dvc(\Hyp')$?
  \espacio{2.8cm}
\end{enumerate}

\nuevapagina
%% ---------------------------------------------------------------
\item \textbf{Leer curvas de aprendizaje.} Para un mismo problema de clasificación se entrenaron dos modelos, A y B, con conjuntos de entrenamiento de distintos tamaños $N$. Las gráficas muestran el error de entrenamiento $\Ein$ y el de validación $E_{\text{val}}$.

\begin{center}
\begin{tikzpicture}
\begin{groupplot}[group style={group size=2 by 1, horizontal sep=1.6cm},
  width=7.6cm, height=5.2cm, xmode=log, ymin=0, ymax=0.45,
  xlabel={$N$}, ylabel={error}, grid=major, grid style={gray!25},
  legend style={font=\small, draw=none, fill=white, fill opacity=0.85, text opacity=1},
  tick label style={font=\small}, label style={font=\small}, title style={font=\bfseries}]
\nextgroupplot[title={Modelo A}, legend pos=north east]
\addplot[thick, blue!70!black, mark=*] coordinates {(100,0.05) (500,0.08) (1000,0.10) (5000,0.12) (10000,0.13)};
\addplot[thick, red!70!black, mark=square*] coordinates {(100,0.40) (500,0.30) (1000,0.25) (5000,0.18) (10000,0.16)};
\legend{$\Ein$, $E_{\text{val}}$}
\nextgroupplot[title={Modelo B}, legend pos=north east]
\addplot[thick, blue!70!black, mark=*] coordinates {(100,0.21) (500,0.22) (1000,0.22) (5000,0.22) (10000,0.22)};
\addplot[thick, red!70!black, mark=square*] coordinates {(100,0.24) (500,0.23) (1000,0.23) (5000,0.23) (10000,0.23)};
\legend{$\Ein$, $E_{\text{val}}$}
\end{groupplot}
\end{tikzpicture}
\end{center}

\begin{enumerate}
  \item ¿Cuál de los dos modelos tiene un problema de sesgo y cuál de varianza? Justifica con lo que ves en las gráficas.
  \espacio{2.3cm}

  \item Con $N = 10{,}000$, ¿cuál usarías? ¿Por qué el $\Ein$ de A \emph{sube} al aumentar $N$?
  \espacio{2.3cm}

  \item Para cada acción, escribe si ayudaría a A, a B, a ambos o a ninguno: (i) conseguir el doble de datos; (ii) aumentar la regularización $\lambda$; (iii) agregar atributos o usar un modelo más flexible.
  \espacio{2.5cm}
\end{enumerate}

\nuevapagina
%% ---------------------------------------------------------------
\item \textbf{La misma prueba en dos fábricas.} Un sistema de visión revisa piezas y marca como defectuosas las que le parecen sospechosas. En el laboratorio se midió que detecta el 90\,\% de las piezas defectuosas ($\text{TVP} = 0.90$) y marca el 5\,\% de las buenas ($\text{TFP} = 0.05$). Se instalará en dos fábricas: en la fábrica 1 el 2\,\% de las piezas son defectuosas; en la fábrica 2, el 20\,\%.
\begin{enumerate}
  \item En cada fábrica, de las piezas marcadas, ¿qué fracción está realmente defectuosa (la precisión)?
  \espacio{3.5cm}

  \item Calcula la exactitud del sistema en la fábrica 1 y compárala con la de un «sistema» que declara buenas todas las piezas. ¿Qué concluyes?
  \espacio{3cm}

  \item TVP y TFP son iguales en las dos fábricas, pero la precisión no. ¿Por qué?
  \espacio{2.5cm}
\end{enumerate}

\end{enumerate}


\end{document}
